\documentclass[usenames,dvipsnames,svgnames,table,aspectratio=169,mathserif]{beamer}

\mode<presentation> {

%\usetheme{default}
\usetheme{Madrid}

% TODO make background not white for BFPG

\setbeamertemplate{footline} % To remove the footer line in all slides uncomment this line

\setbeamertemplate{navigation symbols}{} % To remove the navigation symbols from the bottom of all slides uncomment this line
}
\usepackage{graphicx} % Allows including images
\usepackage{booktabs} % Allows the use of \toprule, \midrule and \bottomrule in tables
\usepackage{hyperref}
\usepackage{apacite}
\usepackage{fancyvrb}
\usepackage{color}
\usepackage{alltt}
\usepackage{listings}
\usepackage{framed}
\usepackage{courier}
\usepackage{minted}
\usepackage{epstopdf}

\hypersetup{colorlinks=false}

\setbeamertemplate{bibliography entry title}{}
\setbeamertemplate{bibliography entry location}{}
\setbeamertemplate{bibliography entry note}{}
\setbeamertemplate{itemize items}[default]
\setbeamertemplate{enumerate items}[default]
\beamertemplatenavigationsymbolsempty
\setbeamertemplate{footline}{}

\newminted{haskell}{}
\newminted{java}{}
\newminted{scala}{}

\definecolor{g}{RGB}{0,100,0}
\newcommand{\highlight}[1]{\colorbox{yellow}{#1}}
\newcommand{\nega}[1]{\colorbox{yellow}{#1}}
\newcommand{\posi}[1]{\colorbox{green}{#1}}
\newcommand{\nl}{\vspace{\baselineskip}}
\newcommand{\pnl}{\pause \nl}

%%----------------------------------------------------------------------------------------
%	TITLE PAGE
%----------------------------------------------------------------------------------------

\title[Type Class: The Ultimate Ad Hoc]{Type Class: The Ultimate Ad Hoc} % The short title appears at the bottom of every slide, the full title is only on the title page
\titlegraphic{\includegraphics[scale=0.2]{data61.eps}}
\author{George Wilson} % Your name
\institute[] % Your institution as it will appear on the bottom of every slide, may be shorthand to save space
{
Data61/CSIRO\\ % Your institution for the title page
\medskip
\href{george.wilson@data61.csiro.au}{george.wilson@data61.csiro.au} % Your email address
}
\date{\today} % Date, can be changed to a custom date

\begin{document}


%%%%%
%%%%% Intro section
%%%%%


\begin{frame}
\titlepage % Print the title page as the first slide
\end{frame}


\begin{frame}

Type classes are a language feature

\begin{itemize}
\item Haskell
\item Purescript
\item Eta
\item Clean
\end{itemize}

or sometimes a design pattern

\begin{itemize}
\item Scala
\end{itemize}

\end{frame}


%%%%%
%%%%% Polymorphism
%%%%%

\begin{frame}
\begin{center}
\huge{Polymorphism}
\end{center}
\end{frame}


\begin{frame}

Something which is {\it polymorphic} has many shapes

\end{frame}


\begin{frame}
\begin{center}
Polymorphism is good


\begin{itemize}
\item less duplication
\item more reuse
\item fewer possible implementations
\end{itemize}
\end{center}
\end{frame}


\begin{frame}
Broadly speaking there are two major forms of polymorphism in programming:

\begin{itemize}
\item {\it parametric} polymorphism
\item {\it ad-hoc} polymorphism
\end{itemize}
\end{frame}


\begin{frame}[fragile]

A {\it parametrically polymorphic} type has at least one {\it type parameter}
which can be instantiated to {\it any type}.

\nl

%Parametrically polymorphic functions behave the same way no matter which type
%they are instantiated to.

%\pnl

Example:

\begin{haskellcode}
reverse :: [a] -> [a]
\end{haskellcode}

\end{frame}


\begin{frame}[fragile]

% TODO make a joke here about "ad-hocly"
An {\it ad-hocly polymorphic} type can be instantiated to some different types,

and may behave differently for each type

\nl

Example:

\begin{javacode}
==
\end{javacode}

\end{frame}


%%%%%
%%%%% INTERFACES SECTION
%%%%%


\begin{frame}
\begin{center}
\huge{Interfaces}
\end{center}
\end{frame}


\begin{frame}[fragile]
\begin{javacode}
interface Equal<A> {
  public boolean eq(A other);
}
\end{javacode}

\pause

\begin{javacode}
class Person implements Equal<Person> {
  public int age;
  public String name;

  public boolean eq(Person other) {
    return this.age == other.age && this.name.equals(other.name);
  }
}
\end{javacode}
\end{frame}


\begin{frame}[fragile]
\begin{javacode}
static <A extends Equal<A>> boolean elementOf(A a, List<A> list) {
  for (A element : list) {
    if (a.eq(element)) return true;
  }
  return false;
}
\end{javacode}

\pnl

\begin{javacode}
elementOf(me, functionalProgrammers);
// true
\end{javacode}
\end{frame}


\begin{frame}[fragile]
\begin{javacode}
package java.lang;

class String {
  private char[] value;
  // other definitions
}
\end{javacode}
\end{frame}


\begin{frame}[fragile]
\begin{javacode}
package java.lang;

class String implements Equal<String> {
  private char[] value;
  // other definitions
}
\end{javacode}
\end{frame}


\begin{frame}[fragile]
\begin{javacode}
class List<A> {
  // implementation details





}
\end{javacode}
\end{frame}


\begin{frame}[fragile]
\begin{javacode}
class List<A> implements Equal<List<A>> {
  // implementation details





}
\end{javacode}
\end{frame}


\begin{frame}[fragile]
\begin{javacode}
class List<A> implements Equal<List<A>> {
  // implementation details

  public boolean eq(List<A> other) {
    // implementation...

  }
}
\end{javacode}
\end{frame}


\begin{frame}[fragile]
\begin{javacode}
class List<A> implements Equal<List<A>> {
  // implementation details

  public boolean eq(List<A> other) {
    // implementation...
    // ... but how do we compare A for equality?
  }
}
\end{javacode}
\end{frame}


\begin{frame}
\begin{itemize}
\item Interface implementation can't be conditional
\item We can only implement interfaces for types we control
\end{itemize}

\end{frame}


%%%%%
%%%%% WHAT ARE TYPE CLASSES
%%%%%


\begin{frame}
\begin{center}
\huge{Type Classes}
\end{center}
\end{frame}


\begin{frame}[fragile]
\begin{haskellcode}
class Equal a where
  eq :: a -> a -> Bool
\end{haskellcode}

\pnl

\begin{haskellcode}
data Person = Person {
  age :: Int
, name :: String
}
\end{haskellcode}

\pnl

\begin{haskellcode}
instance Equal Person where
  eq p1 p2 = eq (age p1) (age p2) && eq (name p1) (name p2)
\end{haskellcode}

\end{frame}


\begin{frame}[fragile]
\begin{haskellcode}
elementOf :: Equal a => a -> [a] -> Bool
elementOf a list =
  case list of
    []    -> False
    (h:t) -> eq a h || elementOf a t
\end{haskellcode}
\end{frame}


\begin{frame}[fragile]
Instances can be constrained

\nl

\begin{haskellcode}
instance (Equal a) => Equal [a] where
  eq []     []     = True
  eq (x:xs) []     = False
  eq []     (y:ys) = False
  eq (x:xs) (y:ys) = eq x y && eq xs ys
\end{haskellcode}

\pnl

We can add type class instances for types we didn't write
\end{frame}


\begin{frame}

\begin{itemize}
\item You can write instances for types you did not write
\item Instances can depend on other instances
\end{itemize}

\end{frame}


\begin{frame}[fragile]

There are exactly two places a type class instance is allowed to exist

\nl

\begin{columns}
\begin{column}[T]{0.49\textwidth}
\begin{block}{\tt Person.hs}
\begin{haskellcode}
data Person = Person
  { age: Int
  , name: String }

instance Equal Person where
  eq p1 p2 = ...
\end{haskellcode}
\end{block}
\end{column}
\begin{column}{0.02\textwidth}
\end{column}
\begin{column}[T]{0.49\textwidth}
\begin{block}{\tt Equal.hs}
\begin{haskellcode}
class Equal a where
  eq :: a -> a -> Bool
\end{haskellcode}
\end{block}
\end{column}
\end{columns}

\end{frame}


\begin{frame}[fragile]

There are exactly two places a type class instance is allowed to exist

\nl

\begin{columns}
\begin{column}[T]{0.49\textwidth}
\begin{block}{\tt Person.hs}
\begin{haskellcode}
data Person = Person
  { age: Int
  , name: String }
\end{haskellcode}
\end{block}
\end{column}
\begin{column}{0.02\textwidth}
\end{column}
\begin{column}[T]{0.49\textwidth}
\begin{block}{\tt Equal.hs}
\begin{haskellcode}
class Equal a where
  eq :: a -> a -> Bool


instance Equal Person where
  eq p1 p2 = ...
\end{haskellcode}
\end{block}
\end{column}
\end{columns}

\end{frame}


\begin{frame}[fragile]

\begin{columns}
\begin{column}[T]{0.49\textwidth}
\begin{block}{\tt Person.hs}
\begin{haskellcode}
data Person = Person
  { age: Int
  , name: String }
\end{haskellcode}
\end{block}
\end{column}
\begin{column}{0.02\textwidth}
\end{column}
\begin{column}[T]{0.49\textwidth}
\begin{block}{\tt Equal.hs}
\begin{haskellcode}
class Equal a where
  eq :: a -> a -> Bool
\end{haskellcode}
\end{block}
\end{column}
\end{columns}

\begin{columns}
\begin{column}{0.5\textwidth}
\begin{block}{\tt EqualInstances.hs}
\begin{haskellcode}
instance Equal Person where
  eq p1 p2 = ...
\end{haskellcode}
\end{block}
\end{column}
\end{columns}
\pause
\begin{center}
``Orphan instance''

Orphan instances break {\it type class coherence}
\end{center}

\end{frame}


\begin{frame}
Type class coherence gives many sane benefits:

\begin{itemize}
\item There can only be zero or one instance 
\item Using an instance never depends on imports
% TODO more here
\end{itemize}
\end{frame}


\begin{frame}
%TODO
TODO show why type classes aren't perfectly flexible

(no custom local instances, maybe show newtypes for sum and product or something)
\end{frame}

%%%%%
%%%%% IMPLICITS (SCALA) SECTION
%%%%%


\begin{frame}
\begin{center}
\Huge{Implicits}

\nl

\large{{\it More Flexible Than Typeclasses}\texttrademark}
\end{center}
\end{frame}


\begin{frame}[fragile]
\begin{scalacode}
case class Person(age: Int, name: String)
\end{scalacode}

\pnl

\begin{scalacode}
trait Equal[A] {
  def eq(a: A, b: A): Boolean
}
\end{scalacode}

\pnl

\begin{scalacode}
implicit def equalPerson: Equal[Person] = new Equal[Person] {
  def eq(a: Person, b: Person): Boolean =
    a.age == b.age && a.name == b.name
}
\end{scalacode}
\end{frame}


\begin{frame}[fragile]
\begin{scalacode}
def elementOf[A](a: A, list: List[A])
                (implicit equalA: Equal[A]): Boolean = {
  list match {
    case Nil => false
    case (h::t) => equal.eq(a, h) || elementOf(a, t)
  }
}
\end{scalacode}
\end{frame}


\begin{frame}[fragile]
\begin{scalacode}
implicit def equalList(implicit equalA: Equal[A]): Equal[List[A]] =
  new Equal[List[A]] {
    def eq(a: List[A], b: List[A]): Boolean = {
      (a,b) match {
        case (Nil,   Nil)   => true
        case (x::xs, Nil)   => false
        case (Nil,   y::ys) => false
        case (x::xs, y::ys) => equalA.eq(x,y) || eq(xs,ys)
      }
    }
  }
\end{scalacode}
\end{frame}


\begin{frame}
\begin{itemize}
\item We can define implicits for types we did not write
\item We can write implicits that depend on implicits
\end{itemize}
\end{frame}


\begin{frame}
\begin{itemize}
\item No restriction on orphan instances
\item No restriction on number of instances
\end{itemize}
\end{frame}


\begin{frame}[fragile]
\begin{scalacode}
sealed trait Ordering
case object LT extends Ordering
case object EQ extends Ordering
case object GT extends Ordering
\end{scalacode}

\pnl

\begin{scalacode}
trait Order[A] {
  def compare(a: A, b: A): Ordering
}
\end{scalacode}

\pnl

\begin{scalacode}
implicit def orderPerson: Order[Person] = new Order[Person] {
  def compare(a: Person, b: Person): Ordering =
    intOrder.compare(a.age, b.age) match {
      case LT => LT
      case EQ => stringOrder.compare(a.name, b.name)
      case GT => GT
    }
}
\end{scalacode}
\end{frame}


\begin{frame}[fragile]
\begin{scalacode}
def sort[A](list: List[A])(implicit orderA: Order[A]): List[A] = {
  // quicksort goes here
}
\end{scalacode}
\end{frame}


\begin{frame}[fragile]
\begin{scalacode}
sort(
  List(
    Person(30, "Robert")
  , Person(20, "John")
  , Person(40, "Alfred")
  )
)
\end{scalacode}

\pnl

\begin{scalacode}
==>

List(
  Person(20, "John")
, Person(30, "Robert")
, Person(40, "Alfred")
)
\end{scalacode}
\end{frame}


\begin{frame}[fragile]
Then the boss says ``I want those sorted by name''.

\pnl

\begin{scalacode}
implicit def orderPersonByName: Order[Person] = new Order[Person] {
  def compare(a: Person, b: Person): Ordering =
    stringOrder.compare(a.name, b.name) match {
      case LT => LT
      case EQ => intOrder.compare(a.age, b.age)
      case GT => GT
    }
}
\end{scalacode}


\end{frame}


\begin{frame}[fragile]
\begin{scalacode}
sort(
  List(
    Person(30, "Robert")
  , Person(20, "John")
  , Person(40, "Alfred")
  )
)
\end{scalacode}

\pnl

\begin{scalacode}
==>

List(
  Person(40, "Alfred")
, Person(20, "John")
, Person(30, "Robert")
)

\end{scalacode}
\end{frame}


\begin{frame}[fragile]
\begin{scalacode}
// both in scope
implicit def orderPerson: Order[Person] = ...
implicit def orderPersonByName: Order[Person] = ...

// what happens?
sort(persons)
\end{scalacode}

\pnl

Hopefully a compiler error!
\end{frame}


% TODO

% show Set example about incoherence doing nasty things quietly
% talk about applying discipline in scala to get coherence


\begin{frame}[fragile]
\begin{block}{Set.scala}
\begin{scalacode}
def emptySet[A]: Set[A]

def insert[A](a: A, set: Set[A])(implicit o: Order[A]): Set[A]

def union[A](s1: Set[A], s2: Set[A])(implicit o: Order[A]): Set[A]
\end{scalacode}
\end{block}
\end{frame}


\begin{frame}
Recommendations when writing implicits:
\begin{itemize}
\item Only create instances in the file that defines the type or the ``type class''
\item Disallow creating more than one instance (regardless of which file you're in)
\end{itemize}

\pnl

What about implicits in external libraries?
\begin{itemize}
\item Assess their usage of implicits. Do they use them as like type classes?
\item If you distrust their implicits, pass everything of theirs explicitly
\end{itemize}
\end{frame}


\begin{frame}

Type classes:
\begin{itemize}
\item Big wins in flexibility, expressiveness, and modularity compared to interfaces
\item Restrictions are straightforward and compiler checked
\item Coherence keeps things sane
\end{itemize}
\end{frame}


\end{document} 

